% problem-000291 12 甲醇合成流程计算
\section*{12. 甲醇合成流程计算}
（图：）

其中，①为甲烷⽓源，压强250.0kPa，温度 $2 5 \ { } ^ { \circ } \mathrm { C } ,$ 流速 $5 5 . 0 \mathrm { m } ^ { 3 } \mathrm { s } ^ { - 1 } ,$ 。②为⽔蒸⽓源，压强200.0kPa，温度 $1 5 0 ~ ^ { \circ } \mathrm { C } ,$ 流速 $1 5 0 . 0 \mathrm { { m } ^ { 3 } \mathrm { { s } ^ { - 1 } } }$ 。合成⽓和剩余反应物的混合物经管路③进⼊ $2 5 ~ ^ { \circ } \mathrm { C }$ 的冷凝器(condenser)，冷凝物由管路⑤流 $\mathbb { H } _ { \circ }$ 。在B中合成的甲醇和剩余反应物的混合物经⑥进⼊ $2 5 ~ ^ { \circ } \mathrm { C }$ 的冷凝器，甲醇冷凝后经管路⑦流出，其密度为 0.791g mL−1。

\subsection*{12-1}
分别写出在步骤A和步骤B中所发⽣的化学反应的⽅程式。

\subsection*{12-2}
假定：所有⽓体皆为理想⽓体，在步骤A和B中完全转化，⽓液在冷凝器中完全分离，计算经步骤A和步骤B后，在⼀秒钟内剩余物的量。

\subsection*{12-3}
实际上，在步骤B中CO的转化率只有三分之⼆。计算在管路⑥中CO、 $\mathrm { H } _ { 2 }$ 和 $\mathrm { C H _ { 3 } O H }$ 的分压（总压强为 10.0 MPa）。

\subsection*{12-4}
当甲醇反应器⾜够⼤，反应达到平衡，管路⑥中的各⽓体的分压服从⽅程：

$
K _ {p} = \frac {p (\mathrm{CH} _ {3} \mathrm{OH}) \times p _ {0} ^ {2}}{p (\mathrm{CO}) \times p ^ {2} (\mathrm{H} _ {2})}
$

式中 ${ } ^ { 1 } p _ { 0 } \ = 0 . 1 0 0 \ \mathrm { M P a }$ ，计算平衡常数 $K _ { p } .$ 。
